\section{The weighted support estimate for grouping}
\label{sec:ope-grouping}
The loss in \eqref{eq:ope-group-weight} pays for two effects. Moving
several marks to a common representative enlarges the supports carrying
them. Expanding the resulting product then creates connected unions
of these enlarged supports. We prove the two counting estimates that
control these operations. They apply to distinct activity labels; no
claim about a logarithmic expansion with repeated labels is needed.

\subsection{Connected unions with distinct labels}
Let $X_i$ be finite supports in a four-dimensional periodic lattice,
indexed by an arbitrary set, and let $c_M(X)$ be their covering number
at a positive integer block size $M$. We use its elementary properties
\[
 |X|\leq(2M+1)^4c_M(X),\qquad
 c_M(X\cup Y)\leq c_M(X)+c_M(Y)\quad\hbox{if }X\cap Y\ne\varnothing.
\]
Let $f_i\geq0$, $a\geq0$, $\gamma\geq0$, and suppose that, for every
site $x$,
\begin{equation}\label{eq:ope-hit-hypothesis}
 \sum_{i:x\in X_i} f_i e^{(a+v)c_M(X_i)}\leq\gamma,
 \qquad (2M+1)^4\gamma\leq v.
\end{equation}
All sums of nonnegative terms have their supremum-of-finite-sums
meaning. A nonempty finite label set $A$ is connected when its
intersection graph, with an edge for overlapping supports, is connected.
Then
\begin{equation}\label{eq:ope-forest-price}
 \sum_{\substack{A\text{ finite, connected}\\x\in\bigcup_{i\in A}X_i}}
 e^{a c_M(\bigcup_{i\in A}X_i)}\prod_{i\in A}f_i
 \leq e^\gamma-1.
\end{equation}

We prove this by a finite recursion. For a finite set of allowed labels
$D$ and a support $Y$, let $\mathcal P(D,Y)$ consist of all subsets
of $D$ whose every connected component has a support meeting $Y$;
the empty subset is included. Set
\[
 F(D,Y)=\sum_{A\in\mathcal P(D,Y)}\prod_{i\in A}p_i,
 \qquad p_i=f_i e^{a c_M(X_i)}.
\]
The following inequality retains each label at most once:
\begin{equation}\label{eq:ope-forest-recursion}
 F(D,Y)\leq\prod_{\substack{i\in D\\X_i\cap Y\ne\varnothing}}
       \bigl(1+p_iF(D\setminus\{i\},X_i)\bigr).
\end{equation}
To see it, choose for each connected component of $A$ one label whose
support meets $Y$. Different components choose different labels.
Deleting the chosen label leaves a set every one of whose connected
components meets its support: follow a path to the deleted vertex
in the old component. Record that remaining set at the chosen label,
and record the empty option at every unchosen label. The union of the
recorded sets with their chosen labels recovers $A$, so the encoding
is injective. Its weight is exactly $\prod_{i\in A}p_i$. Expanding
the product on the right of \eqref{eq:ope-forest-recursion} permits
additional choices, possibly sharing labels between different options;
all their weights are nonnegative. This proves the inequality without
introducing factorial or ordering losses.

Strong induction on $|D|$ gives
\[
 F(D,Y)\leq e^{|Y|\gamma}\leq e^{v c_M(Y)}.
\]
Indeed insert the induction bound in
\eqref{eq:ope-forest-recursion} and use $1+r\leq e^r$:
\[
 F(D,Y)\leq
 \exp\!\left(\sum_{\substack{i\in D\\X_i\cap Y\ne\varnothing}}
           f_i e^{(a+v)c_M(X_i)}\right)
 \leq e^{|Y|\gamma}.
\]
The last inequality sums \eqref{eq:ope-hit-hypothesis} over sites
of $Y$; a support meeting $Y$ is counted at least once. The cardinality
bound and \eqref{eq:ope-hit-hypothesis} prove the second inequality,
which closes the induction. Now take $Y=\{x\}$ and remove the
empty subset, whose weight is one. For a connected $A$, order its
vertices so that each new support meets the preceding union; a
spanning tree supplies such an order. Repeated covering subadditivity
gives $c_M(\bigcup_{i\in A}X_i)\leq\sum_{i\in A}c_M(X_i)$.
Thus each term in \eqref{eq:ope-forest-price} is bounded by its
corresponding product of $p_i$. This proves the desired estimate
for every finite family of connected sets: take $D$ to contain all
their labels. Taking the supremum over finite families proves the
full assertion and its absolute summability.

\subsection{Paying for the new representative sites}
Let a finite set of old marks $I$ be partitioned into groups $V$,
with a chosen old mark in each group and a positive integer bound $h$
on the group size. Write $r_i$ for the old site,
$c_v$ for its group's representative site, and $\pi:I\to V$ for the
partition. Assume that at most $h$ old marks belong to a group and
at most $\ell$ old sites occur at any one lattice site. The same
occupancy bound holds for the representatives, since their selected
old marks are distinct. Suppose
\[
 c_M(\{r_i,c_{\pi(i)}\})\leq k.
\]
An activity indexed by $\alpha$ has a mark set $A_\alpha\subseteq I$
and support $X_\alpha$ containing every $r_i$ for $i\in A_\alpha$.
Enlarge it to
\[
 \widetilde X_\alpha=X_\alpha\cup
              \{c_{\pi(i)}:i\in A_\alpha\}.
\]
Adding the pairs one at a time gives
\begin{equation}\label{eq:ope-enlarged-cover}
 c_M(\widetilde X_\alpha)
       \leq c_M(X_\alpha)+k|A_\alpha|,
\end{equation}
because each pair meets the original support at $r_i$.
For positive mark weights $w_i$ and $0<\eta\leq1$, put
\[
 \widetilde w_i=(\eta/h)w_i e^{-\sigma k},\qquad \sigma\geq0.
\]
For every nonnegative coefficient majorant $q_\alpha$,
\begin{equation}\label{eq:ope-enlargement-price}
 \left(\prod_{i\in A_\alpha}\widetilde w_i\right)q_\alpha
 e^{\sigma c_M(\widetilde X_\alpha)}
 \leq (\eta/h)^{|A_\alpha|}
 \left(\prod_{i\in A_\alpha}w_i\right)q_\alpha
 e^{\sigma c_M(X_\alpha)}.
\end{equation}
This follows immediately from \eqref{eq:ope-enlarged-cover}; the
exponential price paid by each mark cancels its added connecting cover.

We need a local summation version of this observation. Suppose $b,g\geq0$ and
nonnegative quantities $F_\alpha$ have the bound
$\sum_{\alpha:y\in X_\alpha}F_\alpha\leq g$ at every site.
Suppose new weighted quantities $W_\alpha\geq0$ satisfy
$W_\alpha\leq bF_\alpha$ for all $\alpha$, and the stronger
$W_\alpha\leq(b/h)F_\alpha$ whenever $A_\alpha\ne\varnothing$.
Then
\begin{equation}\label{eq:ope-enlarged-hit}
 \sum_{\alpha:y\in\widetilde X_\alpha}W_\alpha
       \leq(1+\ell)bg.
\end{equation}
The old-support contribution is at most $bg$. If $y$ is a new site
of $\widetilde X_\alpha$, some group $v$ has $c_v=y$ and some
$i\in A_\alpha\cap\pi^{-1}(v)$ has $r_i\in X_\alpha$.
Bound this contribution by summing over all such $v$ and all its old
marks. For fixed $i$ the sum is at most $bg/h$; there are at most
$h$ marks in that group, and at most $\ell$ groups at $y$.
This proves \eqref{eq:ope-enlarged-hit}. Activities with no marks
have no added representative sites and are accounted for entirely
in the old-support contribution. This is why the factor $1/h$ is
needed only for activities carrying marks.

\subsection{The resulting local majorant}
Take nonnegative budget constants $T,A,B$. There are three kinds of
entries in the regrouping expansion. The
first are singleton coefficients $u_i$, with $w_i|u_i|\leq T$.
Their local sum, including the support weight, is at most
$\ell T e^\sigma$, because $c_M(\{r_i\})\leq1$ and the occupancy
is at most $\ell$. The second are the nonempty marked coefficients
of the ordinary activity, whose local weighted sum is at most $A$.
Together these have bound
\[
 g_0=\ell T e^\sigma+A.
\]
Both carry a nonempty mark set. Equation
\eqref{eq:ope-enlargement-price} bounds their new weight by
$\eta/h$ times the old weight. Applying
\eqref{eq:ope-enlarged-hit} with $b=\eta$ gives
$(1+\ell)\eta g_0$. Each entry is split into its healthy and unhealthy
support cases by indicator functions. Each piece is bounded by the
same coefficient majorant, so their combined bound is twice this.

The third kind comprises the activity on flagged supports, including
its zero-mark part. Let its original local weighted sum be at most
$B$. For these entries \eqref{eq:ope-enlargement-price} gives a
factor at most one in general, and at most $1/h$ for a nonempty mark
set. Thus \eqref{eq:ope-enlarged-hit} with $b=1$ gives
$(1+\ell)B$. The complete local majorant is therefore
\begin{equation}\label{eq:ope-group-majorant}
 \Gamma=2(1+\ell)\eta(\ell T e^\sigma+A)+(1+\ell)B.
\end{equation}
This bound is independent of the total number of cells and groups.

Apply \eqref{eq:ope-forest-price} with support exponent
$\sigma=\nu+v$, provided $(2M+1)^4\Gamma\leq v$.
Every connected set of distinct entries then has local summed price
at most $e^\Gamma-1$ with exponent $\nu$.
For a component whose mark set is a union of complete groups, let
$W_v=\prod_{i\in\pi^{-1}(v)}\widetilde w_i$.
The product of its old mark weights is exactly the product of these
$W_v$ over its new group marks. Distinct constituent entries have
disjoint marks when their squarefree marker product is nonzero;
there is no multiplicity correction in this equality. Each constituent-label set must yield at most one output entry across
the ordinary and flagged kinds; its support and prescribed flag set
are fixed by its constituents. This uniqueness is proved in
Section~\ref{sec:ope-grouping-algebra}. Grouping the
connected terms by their resulting support and group mark set, the
triangle inequality preserves the local bound $e^\Gamma-1$.
Restrictions to the healthy or prescribed-flag events only reduce
the coefficient majorants. Thus both the ordinary and flagged
connected arrays have this same bound.

For \eqref{eq:ope-group-weight}, the old weights are $\delta W_i$,
$\eta=\delta$, $\sigma=\nu+1/2$, and the maximum group size is $h$.
The new weight of an old mark is precisely
$\delta^2W_i e^{-(\nu+1/2)k}/h$.
Their product over a group is the numerator in
\eqref{eq:ope-group-weight}. The purchased smallness conditions are
$(2M+1)^4\Gamma\leq1/2$ and $e^\Gamma-1\leq1$.
They provide the claimed unit bound for both connected arrays before
the subsequent factor $1/8$ is reserved for their local presentation.
This proves the quantitative support and coefficient estimates used
by regrouping. The exact polynomial decomposition, its gauge/locality
properties and the local projection estimate are separate steps;
these counting bounds do not by themselves establish those identities.
